\begin{helper}
For a sequence of two-bite laws, suppose \(R_n\) holds with high probability,
the event \(\neg G_n\cap R_n\) has negligible probability, and eventually all
sample weights are nonnegative.  Then \(G_n\) holds with high probability.
\end{helper}
\begin{proof}
Let
\[
r_n=\Pr(R_n),\qquad b_n=\Pr(\neg G_n\cap R_n),\qquad g_n=\Pr(G_n).
\]
The hypotheses give \(r_n\to1\) and \(b_n\to0\), hence
\[
r_n-b_n\to1.
\]
For every sufficiently large \(n\), all sample weights are nonnegative.  For
each sample point \(\omega\), the pointwise inequality
\[
1_{R_n}(\omega)w(\omega)-1_{\neg G_n\cap R_n}(\omega)w(\omega)
\le 1_{G_n}(\omega)w(\omega)
\]
holds: if \(G_n(\omega)\) holds it is immediate, and if \(G_n(\omega)\) fails
then the two terms on the left cancel whenever \(R_n(\omega)\) holds and are
both zero otherwise.  Summing over the finite sample space gives
\[
r_n-b_n\le g_n.
\]
The same nonnegativity gives monotonicity of the finite weighted sum, so
\[
g_n\le \Pr(\Omega_n)\le 1,
\]
where the final inequality is
\(\noderef{TwoBiteEventProbabilityTotalMassBound}\).  Therefore \(g_n\) is
squeezed between a sequence converging to \(1\) and the constant sequence
\(1\), so \(g_n\to1\).  This is exactly the asserted
\(\noderef{WithHighProbability}\) conclusion.
\end{proof}
